% 01_introduction.tex
\section{Introduction}
\label{sec:introduction}

The AlphaFold~2 model~\cite{jumper2021highly} and its successor releases
have transformed structural biology by providing high-quality three-dimensional
predictions for virtually every known protein sequence.
The AlphaFold Protein Structure Database~\cite{varadi2022alphafold}
(AlphaFold-DB) currently hosts predictions for over 200 million proteins,
including a complete coverage of the \textit{H.~sapiens} proteome at
database version~4~\cite{tunyasuvunakool2021highly}.
Each predicted structure is accompanied by a per-residue confidence
estimate, the predicted local distance difference test (pLDDT), stored on
a 0--100 scale in the B-factor column of the deposited PDB file.
High pLDDT values ($\gtrsim70$) correlate with well-folded regions;
low values ($\lesssim50$) are associated with intrinsically disordered
regions (IDRs)~\cite{alderson2023systematic}.

Despite this wealth of predictions, systematic proteome-scale analysis of
the topological structure of pLDDT distributions remains limited.
Cazals \& Sarti~\cite{cazals2025alphafold} recently introduced a
rigorous framework, termed Q4 (or Q0 in the preprint's internal
notation), that treats the sequence of pLDDT values along a protein as a
scalar function on a path graph and analyses it through the lens of
topological data analysis~\cite{edelsbrunner2010computational}.
Specifically, residues are inserted into the graph in order of decreasing
pLDDT, and the evolving connected-component structure is tracked via a
Union-Find data structure with the Elder Rule.
The resulting persistence diagram encodes the pLDDT birth-death intervals
of each connected component, providing a compact multi-scale signature of
how confidence varies along the sequence.

From the persistence diagram the authors derive five complementary
statistics.
The fraction of positive-persistence pairs~$f^+_{cp}$ measures the degree
of pLDDT heterogeneity; the normalised Shannon entropy~$H_p$ captures the
diversity of persistence lifetimes; the peak connected-component
count~$N_{cc}^{max}$ reflects the maximum simultaneous fragmentation; and
the count of persistent local maxima~$\mathrm{PLM}(t_p)$ at relative
threshold~$t_p$ identifies the number of topologically significant
confidence peaks.
A high $\mathrm{PLM}$ count is a hallmark of multi-domain proteins, where
distinct structural domains manifest as separate pLDDT peaks separated by
disordered linkers.
The authors demonstrate that $f^+_{cp}$ and $H_p$ are linearly correlated
at Pearson $r \approx 0.97$ across the human proteome, and that applying
the joint filter $n \ge 200$, $H_p \ge 0.25$, $\mathrm{PLM} \ge 3$ at
$t_p = 0.025$ yields approximately 86 candidate multi-domain fragmented
structures in \textit{H.~sapiens}.

\subsection*{Scope of this work}

This report has two goals.
First, we reproduce the Q4 analysis of~\cite{cazals2025alphafold} in full,
starting from raw AlphaFold-DB~v4 PDB files and implementing the
path-graph filtration and all five statistics entirely from scratch in
Python, with only standard scientific libraries
(NumPy~\cite{harris2020numpy},
SciPy~\cite{virtanen2020scipy},
gemmi~\cite{wojdyr2022gemmi}) and
GUDHI~\cite{gudhi2015} for the second-layer PLM persistence step.
We apply the analysis to all 23{,}391 protein fragments in the
\textit{H.~sapiens} proteome that contain at least 200 residues.

Second, we introduce a cross-correlation extension that links the
fragmentation signal from Q4 to the $\mathrm{C}_\alpha$-packing arity
descriptor from Q1~\cite{cazals2025alphafold}.
The arity of a residue is defined as the number of $\mathrm{C}_\alpha$
atoms within 10\,\AA{} of its own $\mathrm{C}_\alpha$; the arity
signature of a protein is the pair $(a_{25}, a_{75})$ of 25th and 75th
percentiles of the per-residue arity distribution.
By computing arity signatures proteome-wide and comparing the arity
distributions of fragmented versus non-fragmented proteins, we test
whether 3D packing density independently predicts the pLDDT fragmentation
phenotype.

\subsection*{Reproducibility and data versions}

Throughout this work we flag a systematic discrepancy that arises from
data-version differences.
AlphaFold-DB~v4 stores pLDDT values with floating-point precision
(a 316-residue protein such as P15121 has 108 distinct pLDDT levels),
whereas the paper appears to have been produced with an earlier version
exhibiting effectively integer-valued pLDDT distributions (typically fewer
than 20 distinct levels per protein).
This difference increases the number of distinct filtration thresholds,
spreading the persistence spectrum and reducing both $H_p$ and the
pairwise Pearson correlation.
Where our results deviate from the paper's targets, we systematically
attribute the deviation to this precision gap and quantify its magnitude.

\subsection*{Organisation}

The remainder of the report is structured as follows.
Section~\ref{sec:background} reviews the necessary background in
path-graph filtrations, persistence diagrams, and the Elder Rule.
Section~\ref{sec:methods} describes the implementation in detail, covering
the filtration algorithm, the five statistics, the PLM computation via
GUDHI, and the arity extension.
Section~\ref{sec:validation} validates the implementation on three
prototype proteins from the paper (ordered P15121, disordered
A0A0G2L439, and mixed Q9VQS4) and on a null random model.
Section~\ref{sec:results} presents the proteome-scale results and the
Q1$\times$Q4 cross-correlation analysis.
Section~\ref{sec:discussion} discusses discrepancies, limitations, and
the broader significance of the arity cross-correlation finding.
Section~\ref{sec:conclusion} concludes.
